% BHU PYQ -- Transcribed & Corrected
% Paper: PHYMV31: Applied Electronics and Fiber Optics
% Exam: B.Sc. (Hons) Semester III (NEP) Examination, 2025-26
% Ref Code: CECONF/N/2025-26/009
% Category: Major / Minor / VAC
% Department: Physics

\documentclass[11pt,a4paper]{article}
\usepackage[a4paper,top=2.5cm,bottom=2.5cm,left=2.8cm,right=2.8cm,headheight=14pt]{geometry}
\usepackage{amsmath,amssymb}
\usepackage[T1]{fontenc}
\usepackage[utf8]{inputenc}
\usepackage{lmodern,microtype,fancyhdr,enumitem,hyperref,lastpage,parskip}

\hypersetup{
    colorlinks=true,
    urlcolor=black,
    linkcolor=black,
    pdftitle={PHYMV31: Applied Electronics and Fiber Optics -- BHU Sem III 2025-26 (NEP)}
}

\pagestyle{fancy}
\fancyhf{}
\renewcommand{\headrulewidth}{0.4pt}
\renewcommand{\footrulewidth}{0.4pt}
\fancyhead[L]{\small \textit{Banaras Hindu University}}
\fancyhead[R]{\small \textit{PHYMV31: Applied Electronics \& Fiber Optics (2025-26)}}
\fancyfoot[C]{\small Page \thepage\ of \pageref{LastPage}}

\newcommand{\pts}[1]{\hfill \textit{[#1]}}

\begin{document}

\begin{center}
  {\large\bfseries BANARAS HINDU UNIVERSITY}\\[2pt]
  {\normalsize B.Sc. (Hons) Semester III Examination, 2025-26 (NEP)}\\[6pt]
  \rule{0.8\linewidth}{0.5pt}\\[6pt]
  {\large\bfseries PHYSICS}\\[4pt]
  {\large\bfseries Paper Code: PHYMV31 : Applied Electronics and Fiber Optics}\\[2pt]
  {\small Ref. No.: CECONF/N/2025-26/009}\\[4pt]
  {\normalsize Time: 2$\frac{1}{2}$ Hours\hfill Maximum Marks: 70}
\end{center}

\rule{\linewidth}{0.4pt}

\smallskip

\noindent \textit{Instructions: Answer five questions including Question 1 which is compulsory. The figures in the right-hand margin indicate marks.}

\bigskip

\noindent \textbf{Question 1.} Answer all questions briefly: \pts{10 $\times$ 2 = 20}
\begin{enumerate}[label=(\alph*)]
  \item What is multimeter? What are its uses?
  \item What is Digital Storage Oscilloscope (DSO)? State one advantage over Cathode Ray Oscilloscope (CRO).
  \item Differentiate between active and passive electronic components.
  \item Describe briefly the depletion layer in a p-n junction diode.
  \item What is field-effect transistor? Mention one of its advantages over a bipolar junction transistor.
  \item A transistor has a current gain $\beta = 100$. Find the value of current gain $\alpha$.
  \item A transformer has 200 turns on its primary coil and 800 turns on its secondary coil. If the primary voltage is 50 V, calculate the secondary voltage of the transformer.
  \item Write two applications of optical fibers.
  \item An optical fiber has a core refractive index $n_1 = 1.500$ and a cladding refractive index $n_2 = 1.414$. Find numerical aperture and acceptance angle in air.
  \item What is a Laser and why is it used in optical fiber communication?
\end{enumerate}

\bigskip

\noindent \textbf{Question 2.}
\begin{enumerate}[label=(\alph*)]
  \item What is a galvanometer? Explain the conversion of a galvanometer into a voltmeter and an ammeter. \pts{6}
  \item A galvanometer has an internal resistance of $50\ \Omega$ and the full-scale deflection current is $2\text{ mA}$. We want to convert it into a 5-volt voltmeter. Find the value of the series resistance. \pts{5+3}
  \item What is the difference between a step-up and a step-down transformer? \pts{4$\frac{1}{2}$}
\end{enumerate}

\bigskip

\noindent \textbf{Question 3.}
\begin{enumerate}[label=(\alph*)]
  \item Explain the construction, working, and applications of a cathode ray oscilloscope (CRO). \pts{6}
  \item With the help of a neat circuit diagram, discuss the working of a Zener diode as a voltage regulator. \pts{6$\frac{1}{2}$}
\end{enumerate}

\bigskip

\noindent \textbf{Question 4.}
\begin{enumerate}[label=(\alph*)]
  \item With the help of a neat circuit diagram, explain the working of a full-wave rectifier. Find the efficiency and the ripple factor of a full-wave rectifier. \pts{7}
  \item Write the current-voltage ($I-V$) equation of a p-n junction diode and obtain the following expression for the dynamic forward resistance $R_f$ of the diode: $R_f = \frac{\eta k T}{e I}$. \pts{5$\frac{1}{2}$}
\end{enumerate}

\bigskip

\noindent \textbf{Question 5.}
\begin{enumerate}[label=(\alph*)]
  \item What is a light-emitting diode? Explain briefly its construction and working. \pts{6}
  \item Draw a neat circuit diagram showing an n-p-n transistor as an amplifier in the common emitter (CE) mode. What is the meaning of the phase shift in the CE mode? \pts{6$\frac{1}{2}$}
\end{enumerate}

\bigskip

\noindent \textbf{Question 6.}
\begin{enumerate}[label=(\alph*)]
  \item Explain the mechanism of total internal reflection (TIR) in optical fibers. Derive the expression for numerical aperture (NA) and discuss how NA influences light-gathering ability. \pts{8}
  \item Explain the terms emitter, base, and collector in the transistors. \pts{4$\frac{1}{2}$}
\end{enumerate}

\bigskip

\noindent \textbf{Question 7.}
\begin{enumerate}[label=(\alph*)]
  \item Describe the different types of optical fiber losses, including absorption, scattering, bending loss and modal dispersion. Explain the physical origin of each type and discuss methods to minimize these losses in practical fiber communication systems. \pts{7}
  \item Explain the difference between the Avalanche breakdown and the Zener breakdown. \pts{5$\frac{1}{2}$}
\end{enumerate}

\bigskip
\begin{center}
  \textbf{***}
\end{center}

\end{document}
